% GENERATED FILE. Edit canonical lessons, then run npm run generate:books.
% canonical-source-hash: fnv1a64:3e8c42bf40a8e226
% canonical-lessons: AR-C498-asada, AR-C498-ahzana, AR-C498-adhaka, AR-C498-aghdaba, AR-C498-lama

\chapter{Making others feel, and blame}
\label{ch:ar-making-others-feel-and-blame}
\hlchaptermodality{498}

\begin{chapteropening}
\textbf{By the end of this chapter:} \emph{I can say that something made someone happy, sad, laugh or angry, and that someone blamed someone.}

Complete the last lesson of chapter 498: I can say that something made someone happy, sad, laugh or angry, and that someone blamed someone.
\end{chapteropening}

\section[\texorpdfstring{\ar{أسعد} (asʿada)}{أسعد (asʿada)}]{\ar{أسعد} (asʿada) — to make happy}

\label{lesson:AR-C498-asada}



\begin{quote}
Before the new one: say the Arabic for hesitant, then the Arabic for to hold back.
\end{quote}

\subsection*{You'll want to know: \ar{أسعد}}
\textbf{\ar{أسعد}} — \emph{asʿada} — \textquotedblleft{}to make happy\textquotedblright{}.

\begin{quote}
\ar{أسعدت} \ar{الهدية} \ar{جدتي} \ar{كثيرا}.
\end{quote}

\begin{quote}
\emph{asʿadat al-hadiyya jaddatī kathīran.}
\end{quote}

\begin{quote}
The present made my grandmother very happy.
\end{quote}

Present: \ar{يسعد}, \emph{yusʿid}.

\begin{cousinweb}
\textbf{\ar{أسعد}} (\emph{asʿada}) is built on the native Arabic root \emph{s-a-d}, \textquotedblleft{}to be lucky, to be happy\textquotedblright{} (Lane). It is Form IV: to make someone saʿīd.
\end{cousinweb}

\begin{grammarlens}[title={What you've built}]
One more word for describing what you lived and felt.
\end{grammarlens}

\subsection*{Your turn}
\begin{itemize}
\raggedright
  \item \emph{Say it:} \emph{asʿada}
  \item \emph{Say it:} \emph{asʿada}, once more
  \item \emph{Say it:} \emph{asʿadat al-hadiyya jaddatī kathīran.} — The present made my grandmother very happy.
  \item \emph{Recall:} say the Arabic for hesitant, then the Arabic for to hold back, then say \emph{asʿada} again
\end{itemize}

\begin{tcolorbox}[breakable,colback=teal!4,colframe=teal!35!black,title={Before you move on}]
What does \ar{أسعد} mean? (\textquotedblleft{}To make happy\textquotedblright{}.) Say it once more.
\end{tcolorbox}
\section[\texorpdfstring{\ar{أحزن} (aḥzana)}{أحزن (aḥzana)}]{\ar{أحزن} (aḥzana) — to sadden}

\label{lesson:AR-C498-ahzana}



\begin{quote}
Before the new one: say the Arabic for to hold back, then the Arabic for to make happy.
\end{quote}

\subsection*{You'll want to know: \ar{أحزن}}
\textbf{\ar{أحزن}} — \emph{aḥzana} — \textquotedblleft{}to sadden\textquotedblright{}.

\begin{quote}
\ar{أحزن} \ar{الفيلم} \ar{كل} \ar{الناس}.
\end{quote}

\begin{quote}
\emph{aḥzana al-fīlm kull an-nās.}
\end{quote}

\begin{quote}
The film saddened everyone.
\end{quote}

Present: \ar{يحزن}, \emph{yuḥzin}.

Read the present \ar{يحزن} as yuḥzin here, yaḥzan for to be sad.

\begin{cousinweb}
\textbf{\ar{أحزن}} (\emph{aḥzana}) is built on the native Arabic root \emph{ḥ-z-n}, \textquotedblleft{}grief, sadness\textquotedblright{} (Lane). It is Form IV: to cause grief.
\end{cousinweb}

\begin{grammarlens}[title={What you've built}]
One more word for describing what you lived and felt.
\end{grammarlens}

\subsection*{Your turn}
\begin{itemize}
\raggedright
  \item \emph{Say it:} \emph{aḥzana}
  \item \emph{Say it:} \emph{aḥzana}, once more
  \item \emph{Say it:} \emph{aḥzana al-fīlm kull an-nās.} — The film saddened everyone.
  \item \emph{Recall:} say the Arabic for to hold back, then the Arabic for to make happy, then say \emph{aḥzana} again
\end{itemize}

\begin{tcolorbox}[breakable,colback=teal!4,colframe=teal!35!black,title={Before you move on}]
What does \ar{أحزن} mean? (\textquotedblleft{}To sadden\textquotedblright{}.) Say it once more.
\end{tcolorbox}
\section[\texorpdfstring{\ar{أضحك} (aḍḥaka)}{أضحك (aḍḥaka)}]{\ar{أضحك} (aḍḥaka) — to make laugh}

\label{lesson:AR-C498-adhaka}



\begin{quote}
Before the new one: say the Arabic for to make happy, then the Arabic for to sadden.
\end{quote}

\subsection*{You'll want to know: \ar{أضحك}}
\textbf{\ar{أضحك}} — \emph{aḍḥaka} — \textquotedblleft{}to make laugh\textquotedblright{}.

\begin{quote}
\ar{أضحكتنا} \ar{قصة} \ar{جدي}.
\end{quote}

\begin{quote}
\emph{aḍḥakatnā qiṣṣat jaddī.}
\end{quote}

\begin{quote}
My grandfather's story made us laugh.
\end{quote}

Present: \ar{يضحك}, \emph{yuḍḥik}.

Muḍḥik, funny, is its participle.

\begin{cousinweb}
\textbf{\ar{أضحك}} (\emph{aḍḥaka}) is built on the native Arabic root \emph{ḍ-ḥ-k}, \textquotedblleft{}to laugh\textquotedblright{} (Lane). It is Form IV of ḍaḥika, to laugh: to make someone laugh.
\end{cousinweb}

\begin{grammarlens}[title={What you've built}]
One more word for describing what you lived and felt.
\end{grammarlens}

\subsection*{Your turn}
\begin{itemize}
\raggedright
  \item \emph{Say it:} \emph{aḍḥaka}
  \item \emph{Say it:} \emph{aḍḥaka}, once more
  \item \emph{Say it:} \emph{aḍḥakatnā qiṣṣat jaddī.} — My grandfather's story made us laugh.
  \item \emph{Recall:} say the Arabic for to make happy, then the Arabic for to sadden, then say \emph{aḍḥaka} again
\end{itemize}

\begin{tcolorbox}[breakable,colback=teal!4,colframe=teal!35!black,title={Before you move on}]
What does \ar{أضحك} mean? (\textquotedblleft{}To make laugh\textquotedblright{}.) Say it once more.
\end{tcolorbox}
\section[\texorpdfstring{\ar{أغضب} (aghḍaba)}{أغضب (aghḍaba)}]{\ar{أغضب} (aghḍaba) — to make angry}

\label{lesson:AR-C498-aghdaba}



\begin{quote}
Before the new one: say the Arabic for to sadden, then the Arabic for to make laugh.
\end{quote}

\subsection*{You'll want to know: \ar{أغضب}}
\textbf{\ar{أغضب}} — \emph{aghḍaba} — \textquotedblleft{}to make angry\textquotedblright{}.

\begin{quote}
\ar{أغضب} \ar{التأخير} \ar{المدير}.
\end{quote}

\begin{quote}
\emph{aghḍaba at-taʾkhīr al-mudīr.}
\end{quote}

\begin{quote}
The delay made the manager angry.
\end{quote}

Present: \ar{يغضب}, \emph{yughḍib}.

\begin{cousinweb}
\textbf{\ar{أغضب}} (\emph{aghḍaba}) is built on the native Arabic root \emph{gh-ḍ-b}, \textquotedblleft{}anger, wrath\textquotedblright{} (Lane). It is Form IV: to cause anger.
\end{cousinweb}

\begin{grammarlens}[title={What you've built}]
One more word for describing what you lived and felt.
\end{grammarlens}

\subsection*{Your turn}
\begin{itemize}
\raggedright
  \item \emph{Say it:} \emph{aghḍaba}
  \item \emph{Say it:} \emph{aghḍaba}, once more
  \item \emph{Say it:} \emph{aghḍaba at-taʾkhīr al-mudīr.} — The delay made the manager angry.
  \item \emph{Recall:} say the Arabic for to sadden, then the Arabic for to make laugh, then say \emph{aghḍaba} again
\end{itemize}

\begin{tcolorbox}[breakable,colback=teal!4,colframe=teal!35!black,title={Before you move on}]
What does \ar{أغضب} mean? (\textquotedblleft{}To make angry\textquotedblright{}.) Say it once more.
\end{tcolorbox}
\section[\texorpdfstring{\ar{لام} (lāma)}{لام (lāma)}]{\ar{لام} (lāma) — to blame}

\label{lesson:AR-C498-lama}



\begin{quote}
Before the new one: say the Arabic for to make laugh, then the Arabic for to make angry.
\end{quote}

\subsection*{You'll want to know: \ar{لام}}
\textbf{\ar{لام}} — \emph{lāma} — \textquotedblleft{}to blame\textquotedblright{}.

\begin{quote}
\ar{لامني} \ar{المدير} \ar{على} \ar{التأخير}.
\end{quote}

\begin{quote}
\emph{lāmanī al-mudīr ʿalā at-taʾkhīr.}
\end{quote}

\begin{quote}
The manager blamed me for the delay.
\end{quote}

Present: \ar{يلوم}, \emph{yalūm}.

It takes \ar{على} for the fault.

\begin{cousinweb}
\textbf{\ar{لام}} (\emph{lāma}) is built on the native Arabic root \emph{l-w-m}, \textquotedblleft{}to blame, to censure\textquotedblright{} (Lane). Its middle root letter w shows as alif in the past and as ū in the present.
\end{cousinweb}

\begin{grammarlens}[title={What you've built}]
One more word for describing what you lived and felt.
\end{grammarlens}

\subsection*{Your turn}
\begin{itemize}
\raggedright
  \item \emph{Say it:} \emph{lāma}
  \item \emph{Say it:} \emph{lāma}, once more
  \item \emph{Say it:} \emph{lāmanī al-mudīr ʿalā at-taʾkhīr.} — The manager blamed me for the delay.
  \item \emph{Recall:} say the Arabic for to make laugh, then the Arabic for to make angry, then say \emph{lāma} again
\end{itemize}

\begin{tcolorbox}[breakable,colback=teal!4,colframe=teal!35!black,title={Before you move on}]
What does \ar{لام} mean? (\textquotedblleft{}To blame\textquotedblright{}.) Say it once more.
\end{tcolorbox}
